\section{分次子代数,分次代数的分次理想}

记号约定:$A$是$\Delta$型分次环,$E$是$\Delta$型分次$A$-代数.

\begin{definition}{$\Delta$型分次$A$-子代数}{}
	设$F$是$E$的子代数.如果$F$是$E$的分次$A$-子模,则称$F$是$E$的$\Delta$型分次($A$-)子代数.
\end{definition}

\begin{remark}{}{}
	记$\left(E_{\lambda}\right)_{\lambda\in\Delta}$为$E$的分次,则$F$是$E$的分次$A$-子模表明$\left(F\cap E_{\lambda}\right)_{\lambda\in\Delta}$是$F$的分次,典范单射$F\to E$是$\Delta$型分次代数同态.
\end{remark}

\begin{definition}{$\Delta$型分次(左,右,双边)理想}{}
	\begin{enumerate}
		\item 设$\mathfrak{a}$是$E$的左(右)理想.若$\mathfrak{a}$是$E$的$\Delta$型分次左(右)$A$-子模,则称$\mathfrak{a}$是$E$的$\Delta$型分次左(右)理想.
		\item 设$\mathfrak{a}$是$E$的双边理想.若$\mathfrak{a}$是$E$的$\Delta$型分次左代数理想和右代数理想,则称$\mathfrak{a}$是$E$的$\Delta$型分次双边理想.
	\end{enumerate}
\end{definition}

\begin{proposition}{商代数的分次结构}{}
	若$\mathfrak{a}$是$E$的$\Delta$型分次双边代数理想,则$E/\mathfrak{a}$的商分次与$E/\mathfrak{a}$的$A$-代数结构相容,典范映射$E\to E/\mathfrak{a}$是分次代数同态.
\end{proposition}

\begin{proposition}{分次代数同态的像与核}{}
	设$F$是$\Delta$型分次代数,$u:E\to F$是$\Delta$型分次代数同态.
	\begin{enumerate}
		\item $\im\left(u\right)$是$F$的$\Delta$型分次子代数,$\ker\left(u\right)$是$U$的$\Delta$型分次双边理想.
		\item $u$诱导的典范双射$E/\ker\left(u\right)\to\im\left(u\right)$是$\Delta$型分次$A$-代数同构.
	\end{enumerate}
\end{proposition}

\begin{proposition}{齐次元生成的代数结构}{}
	进一步假设$A$是交换环,$S$是$E$的一些齐次元构成的集合.
	\begin{enumerate}
		\item $S$生成的$A$-子代数是$\Delta$型分次$A$-子代数.
		\item $S$生成的左(右,双边)理想都是$\Delta$型分次左(右,双边)理想.
	\end{enumerate}
\end{proposition}





